%% =============================================================
%%  chapters/02_background.tex
%% =============================================================
\chapter{Background and Fundamentals}
\label{ch:background}

\section{Cabin Monitoring Systems}
\label{sec:cabin-monitoring-systems}
Cabin monitoring systems are commonly divided into driver monitoring
systems (DMS), which focus on the driver's attentiveness and drowsiness,
and occupant monitoring systems (OMS), which assess the presence,
position, and classification of all vehicle occupants. Regulatory bodies
such as Euro NCAP have progressively increased the weight given to
occupant status monitoring in their safety rating protocols, driving
adoption of both camera-based and pressure-based sensing across vehicle
segments \cite{euroncap2023protocol}. Typical OMS use cases include
seatbelt reminder logic, smart airbag suppression for child seats, and
post-crash occupant status reporting for emergency services.

%% TODO: Add a timeline of relevant Euro NCAP / GTR regulation milestones.

\section{Pressure Sensing Technology}
\label{sec:pressure-sensing}

\subsection{Working Principle and Sensor Types}
\label{sec:pressure-working-principle}
Seat-integrated pressure sensors typically use resistive, capacitive, or
piezoelectric transduction principles arranged in a grid to produce a
spatially resolved pressure map of the seat surface. Resistive force
sensing resistor (FSR) arrays are low-cost and widely used in automotive
seat occupancy classification, while capacitive mats offer higher
sensitivity at increased manufacturing cost \cite{mueller2019pressure}.
The choice of sensor technology directly affects spatial resolution,
sampling rate, and susceptibility to drift over the vehicle's operating
temperature range.

\cite{s21103346} This is paper on lower resolution pressure sensor array for driver posture monitoring and pose estimation.

%%
%% TODO: Document the sensor mat setup and grid resolution used.


\subsection{Pressure Heatmap Representation}
\label{sec:pressure-heatmap-representation}
The raw output of a pressure sensor grid is commonly represented as a
two-dimensional heatmap, where each pixel encodes the normalised pressure
value at a corresponding physical location on the seat surface. This
representation allows pressure data to be processed with the same
convolutional architectures typically applied to single-channel images,
which is a key enabler for the fusion architecture proposed in
Chapter~\ref{ch:fusion-architecture}.

%% TODO: Include the actual sensor grid resolution (rows x columns) used.

\section{Camera Systems for In-Cabin Use}
\label{sec:camera-systems}

\subsection{RGB and IR Cameras}
\label{sec:rgb-ir-cameras}
Standard RGB cameras provide high-resolution colour imagery suitable for
detailed pose and appearance analysis but perform poorly under low-light
or backlit conditions common in vehicle cabins. Near-infrared (IR)
cameras paired with active IR illumination maintain consistent image
quality irrespective of ambient lighting and are therefore preferred in
production driver and occupant monitoring systems \cite{schmidt2020driver}.

%% TODO: Specify camera sensor model, resolution, and frame rate used.

\subsection{Depth Cameras }
\label{sec:depth-tof-cameras}

%% TODO: State whether a depth channel is used in this thesis or reserved
%% for future work.

\section{Sensor Fusion Fundamentals}
\label{sec:sensor-fusion-fundamentals}

\subsection{Fusion Levels}
\label{sec:fusion-levels}
Sensor fusion approaches are typically categorised into data-level fusion
(combining raw signals before any feature extraction), feature-level
fusion (combining learned intermediate representations), and
decision-level fusion (combining independent per-modality predictions).
Feature-level fusion generally offers the best balance between
representational flexibility and robustness to modality-specific noise,
and is the strategy adopted for the architecture in this thesis
\cite{chen2020multimodal}.




%% TODO: 
%% in Chapter 8, or whether this section remains purely for background.

\subsection{Deep Learning-Based Fusion}
\label{sec:deep-learning-fusion}
Deep learning-based fusion replaces hand-crafted process and measurement
models with learned convolutional or attention-based feature extractors
per modality, followed by a learned fusion module, building on the
general deep learning foundations of representation learning and
gradient-based training described in \cite{goodfellow2016deep}. This
approach scales naturally to high-dimensional, spatially structured
inputs such as camera images and pressure heatmaps, and can be trained
end-to-end jointly with the downstream task heads
\cite{oliveira2022crossmodal}.

\section{Multi-Task Learning}
\label{sec:multi-task-learning}
Multi-task learning trains a single shared backbone with multiple
task-specific output heads, allowing related tasks to regularise one
another and share statistical strength, which is particularly valuable
when labelled data for some tasks (e.g. weight estimation) is scarcer than
for others (e.g. binary occupancy). A central challenge in multi-task
learning is balancing the loss contributions of tasks with different
scales and convergence rates, which is addressed explicitly for this
thesis in Section~\ref{sec:loss-function-design}.
%% TODO: Cite foundational multi-task learning literature (e.g. Caruana,
%% Kendall et al. uncertainty weighting) once the reference list is final.

\section{Human Body Modelling for Simulation}
\label{sec:human-body-modelling}
Parametric human body models represent body shape and pose using a
low-dimensional set of shape and joint-angle parameters, enabling
generation of anatomically plausible body meshes across a wide population
of statures and weights. Combined with a physics-based seat contact
simulation, such models allow synthetic pressure heatmaps to be rendered
for arbitrary occupant configurations without requiring a physical
subject, which forms the basis of the synthetic data generation pipeline
described in Section~\ref{sec:synthetic-heatmap-generation}
\cite{ivanov2021synthetic}.
%% TODO: Name the specific body model (e.g. SMPL or similar) once selected.
